% GENERATED FILE. Edit canonical lessons, then run npm run generate:books.
% canonical-source-hash: fnv1a64:937ff4abf5894ba8
% canonical-lessons: ES-C467-producir, ES-C467-suficiente, ES-C467-alimentar, ES-C467-objetivo, ES-R467-repaso-alimento, ES-C467-sintesis-alimento

\chapter{El huerto comunitario}
\label{ch:es-alimento}
\hlchaptermodality{467}

\begin{chapteropening}
\textbf{By the end of this chapter:} \emph{I can read a monthly production sheet on a community garden's door: understand what was produced, who it fed, what the target was, and why a notice writes 'no fue suficiente' where a person would say 'no fue bastante'.}

Read a community garden's monthly figures, then say whether the target was met.
\end{chapteropening}

\section[producir]{producir — to produce}

\label{lesson:ES-C467-producir}



\begin{quote}
Say \textbf{reducir} and \textbf{la ducha}. The verb coming next is their close relative, and it will not make you learn anything twice.
\end{quote}

\subsection*{You'll want to know: producir}
\textbf{producir} — to produce, to make.

\begin{quote}
El huerto \textbf{produce} unos cien kilos al mes.
\end{quote}

\emph{Producir comida} — to produce food. \emph{La producción} is the noun, and \textbf{el producto} is the thing that comes out.

\begin{grammarlens}[title={Grammar Lens: nothing new here, which is the point}]
\emph{Reducir} taught you this whole verb, five chapters ago. \emph{Producir} is the same class and takes the same two irregularities:

\noindent
\begin{tabularx}{\linewidth}{@{}>{\raggedright\arraybackslash}X>{\raggedright\arraybackslash}X>{\raggedright\arraybackslash}X@{}}
\toprule
\textbf{where} & \textbf{reducir} & \textbf{producir} \\
\midrule
the \textbf{z} before \emph{c} & redu\textbf{zc}o & produ\textbf{zc}o \\
the \textbf{j} past & redu\textbf{j}e & produ\textbf{j}e \\
\bottomrule
\end{tabularx}

So the subjunctive is \emph{produzca} — the form on a notice, as \emph{reduzca} was on the road sign — and the past is \emph{produjo}, never \emph{producí}.

That is worth stopping on. Learning \emph{reducir} bought you the whole \emph{-ducir} corner of that class — \emph{producir}, and \emph{traducir} and \emph{introducir} when you meet them — the verbs that take the \textbf{j} past as well as the \textbf{z}.

Not the whole class, though. \emph{Reducir}'s own lesson warned you: \emph{conocer} shares the \textbf{-zc-} and not the past, going \emph{conocí}. Two groups, one spelling, and the \emph{-ducir} verbs are the half that also changes its past.
\end{grammarlens}

\begin{cousinweb}
\textbf{pro-}, forward, plus \textbf{ducere}, to lead — the \emph{ducere} already under \emph{reducir} and under \emph{la ducha}.

To produce is to \textbf{lead something forward}, out where it can be seen. A farmer leads a crop forward out of the ground; a factory leads goods forward out of a door. English keeps the picture in \textbf{produce} as a noun, the stuff itself, and in a theatre \textbf{producer}, who leads the show out in front of an audience.

Set them side by side and the prefix does all the work: \emph{reducir} leads back, \emph{producir} leads forward, and \emph{conducir} — which \emph{reducir}'s table already showed you — leads along.
\end{cousinweb}

\subsection*{Your turn}
Say these aloud:

\begin{itemize}
\raggedright
  \item the garden produces about a hundred kilos a month
  \item I produce, using the first-person form
  \item the past of producir, the one that is not producí
\end{itemize}

\begin{tcolorbox}[breakable,colback=teal!4,colframe=teal!35!black,title={Before you move on}]
To produce? (\textbf{Producir}.) I produce? (\textbf{Produzco}.) It produced? (\textbf{Produjo}.) Which way does \emph{pro-} lead? (\textbf{Forward}.)
\end{tcolorbox}
\section[suficiente]{suficiente — enough, and only enough}

\label{lesson:ES-C467-suficiente}



\begin{quote}
Say \textbf{bastante} and \textbf{escaso}. You already own both ends of this. The word coming next overlaps one of them, and the overlap is the lesson.
\end{quote}

\subsection*{You'll want to know: suficiente}
\textbf{suficiente} — enough, sufficient.

\begin{quote}
No producimos comida \textbf{suficiente} para todos.
\end{quote}

\emph{Comida suficiente} — enough food. \emph{Lo suficiente} — enough of it, as a noun. \emph{Suficientemente} — sufficiently.

Like \emph{bastante}, it has one form for both genders: \emph{un plato suficiente}, \emph{una cantidad suficiente}.

Said of a person it turns unkind — \emph{un tono suficiente} is a smug one.

\begin{grammarlens}[title={Grammar Lens: bastante already means this, and that is the problem}]
\emph{Bastante} was given to you as the middle of a scale — \emph{poco}, \textbf{bastante}, \emph{muy} — and its lesson was firm that it does \textbf{not} mean \emph{a lot}. So what is \emph{suficiente} for?

\noindent
\begin{tabularx}{\linewidth}{@{}>{\raggedright\arraybackslash}X>{\raggedright\arraybackslash}X@{}}
\toprule
\textbf{word} & \textbf{what else it does} \\
\midrule
bastante & \emph{quite}, \emph{fairly} — \emph{bastante caro} \\
suficiente & names a requirement, not a place on a scale \\
\bottomrule
\end{tabularx}

\emph{Bastante} has that second, adverbial job: \emph{bastante caro} is \textbf{fairly} expensive, a reading on the scale. \emph{Suficiente} has no scale reading at all. It measures a thing against what was needed, and reports whether the need was met.

Both can say there is not enough — \emph{no hay bastante} is the everyday phrasing. A notice reaches for \emph{no hay suficiente} because it names the requirement rather than the quantity, and because it is the formal choice.
\end{grammarlens}

\begin{cousinweb}
Latin \textbf{sufficere}: \textbf{sub-}, up from under, plus \textbf{facere}, to make or do — the \emph{facere} that \emph{la factura} called one of the hardest-working roots in the language.

To \emph{sufficere} was to put something \textbf{under} a thing: supply it from below, prop it, make it stand. Adequacy came out of that, which is why sufficiency is load-bearing rather than abundant — the amount that keeps the thing up.

\emph{Suficiente} wears the \textbf{-ente} of \emph{exigente}; \emph{bastante} wears \textbf{-ante}. They are the same ending choosing its vowel by verb family — \emph{-ar} verbs take \emph{-ante}, \emph{-er} and \emph{-ir} verbs take \emph{-ente}.

And the two words are one language at two speeds, the pattern you have from \emph{rueda} against \emph{rotativo}. \emph{Suficiente} is a real Latin present participle, \emph{sufficientem}, lifted off the page. \emph{Bastante} was built inside Spanish on the inherited verb \emph{bastar} — there was never a Latin \emph{bastans}.
\end{cousinweb}

\subsection*{Your turn}
Say these aloud:

\begin{itemize}
\raggedright
  \item we do not produce enough food for everybody
  \item which of the two can also mean \textquotedblleft{}fairly\textquotedblright{}
  \item what sub- plus facere does under a thing
\end{itemize}

\begin{tcolorbox}[breakable,colback=teal!4,colframe=teal!35!black,title={Before you move on}]
Enough, on a notice? (\textbf{Suficiente}.) Which word also means \emph{fairly}? (\textbf{Bastante}.) What does \emph{sufficere} do? (\textbf{Props a thing from underneath} — adequacy is load-bearing.)
\end{tcolorbox}
\section[alimentar]{alimentar — to feed}

\label{lesson:ES-C467-alimentar}



\begin{quote}
Say \textbf{el medicamento} and \textbf{producir}. One lends its ending to this chapter's noun; the other is what has to happen before you can do the verb.
\end{quote}

\subsection*{You'll want to know: alimentar}
\textbf{alimentar} — to feed, to nourish.

\begin{quote}
Lo que producimos \textbf{alimenta} a unas cuarenta familias.
\end{quote}

\emph{Alimentar a un niño} — to feed a child; the \textbf{a} is there because the thing being fed is a person. \textbf{El alimento} is the food itself, and \textbf{la alimentación} is the feeding, or a person's diet.

\emph{Alimentarse de algo} is what an animal or a person lives on: \emph{se alimentan de fruta}.

\begin{cousinweb}
Latin \textbf{alere} was wider than \emph{to feed}. It meant to nourish, and by that to \textbf{rear} — to bring a thing on, to make it bigger. Feeding is the sense that stuck in Spanish; the others are worth having:

\noindent
\begin{tabularx}{\linewidth}{@{}>{\raggedright\arraybackslash}X>{\raggedright\arraybackslash}X@{}}
\toprule
\textbf{English} & \textbf{how the root shows} \\
\midrule
\textbf{adult} & somebody who has finished growing \\
\textbf{alumnus} & somebody being brought up \\
\bottomrule
\end{tabularx}

An \textbf{alumnus} is literally \emph{one being nourished} — a foster-child, then a pupil, then a graduate. \textbf{Coalesce} is \emph{co-} plus \emph{alescere}, to grow together.

So \emph{alimentar} is not about hunger. It is about making something bigger, and food is the means.
\end{cousinweb}

\begin{grammarlens}[title={Grammar Lens: -mento, and the noun it builds here}]
\emph{El medicamento} gave you \textbf{-mento}, the formal twin of \emph{-miento}, arriving by the other road. Latin \emph{-mentum} often named the \textbf{means} of a thing, and both of these are that kind:

\noindent
\begin{tabularx}{\linewidth}{@{}>{\raggedright\arraybackslash}X>{\raggedright\arraybackslash}X@{}}
\toprule
\textbf{noun} & \textbf{the means of} \\
\midrule
el medica\textbf{mento} & healing \\
el ali\textbf{mento} & growing \\
\bottomrule
\end{tabularx}

Watch the cut: it is \emph{ali-} plus \emph{-mento}, exactly as \emph{medica-} plus \emph{-mento}. And the noun is older than the verb — \textbf{Latin} built \emph{alimentum} off \emph{alere} first, and \emph{alimentare} came afterwards. Spanish inherited both, which is why the noun is \emph{alimento} and not \emph{alimentamento}.
\end{grammarlens}

\subsection*{Your turn}
Say these aloud:

\begin{itemize}
\raggedright
  \item what we produce feeds about forty families
  \item the food itself, then the diet
  \item what alere meant before it meant feeding
\end{itemize}

\begin{tcolorbox}[breakable,colback=teal!4,colframe=teal!35!black,title={Before you move on}]
To feed? (\textbf{Alimentar}.) The food? (\textbf{El alimento}.) What did \emph{alere} cover? (\textbf{Nourishing, and rearing} — making a thing bigger.)
\end{tcolorbox}
\section[el objetivo]{el objetivo — the target}

\label{lesson:ES-C467-objetivo}



\begin{quote}
Say \textbf{el objeto} and \textbf{rotativo}. One is the word underneath the next; the other is its ending.
\end{quote}

\subsection*{You'll want to know: el objetivo}
\textbf{el objetivo} — the target, the aim, the objective.

\begin{quote}
\textbf{Objetivo} del mes: alimentar a cincuenta familias.
\end{quote}

\emph{Cumplir el objetivo} — to meet the target. \emph{Nuestro objetivo es…} — our aim is… As an adjective it means \textbf{objective}: \emph{un informe objetivo}.

\begin{cousinweb}
You have the whole of this already. \emph{El objeto} was \textbf{ob-}, before or against, plus \textbf{iacere}, to throw — a thing \emph{thrown before} you, put in your way. That lesson listed the family: \emph{object}, \emph{project}, \emph{inject}, \emph{eject}, \emph{subject}, all of them throws.

\emph{Objetivo} is that same thrown thing wearing \textbf{-ivo}: \emph{to do with the thing placed in front}. The throw does not turn round — the grammatical object is already what an action is thrown at, as \emph{el objeto}'s own chapter said. What narrowed is \textbf{why} the thing is there.

\noindent
\begin{tabularx}{\linewidth}{@{}>{\raggedright\arraybackslash}X>{\raggedright\arraybackslash}X@{}}
\toprule
\textbf{word} & \textbf{the thing placed before you} \\
\midrule
el objeto & whatever is there \\
el objetivo & the one put there to be reached \\
\bottomrule
\end{tabularx}

English kept both jobs on the same word: an \emph{objective} is an aim, and an \textbf{objective lens} is the one that faces the object.
\end{cousinweb}

\begin{grammarlens}[title={Grammar Lens: -ivo makes an adjective that can stand alone}]
\emph{Rotativo} gave you \textbf{-ivo}: an ending that names a habit or a capacity — of the kind that does this.

What \emph{objetivo} adds is that such an adjective can harden into a noun and keep both jobs:

\noindent
\begin{tabularx}{\linewidth}{@{}>{\raggedright\arraybackslash}X>{\raggedright\arraybackslash}X@{}}
\toprule
\textbf{use} & \textbf{example} \\
\midrule
adjective & \emph{un informe objetivo} \\
noun & \emph{el objetivo del mes} \\
\bottomrule
\end{tabularx}

Spanish does this often and English follows it here, which makes the pair easy to trust: \emph{an objective report}, \emph{the month's objective}.
\end{grammarlens}

\subsection*{Your turn}
Say these aloud:

\begin{itemize}
\raggedright
  \item this month's target: to feed fifty families
  \item an objective report
  \item the difference between an objeto and an objetivo
\end{itemize}

\begin{tcolorbox}[breakable,colback=teal!4,colframe=teal!35!black,title={Before you move on}]
The target? (\textbf{El objetivo}.) What is inside it? (\emph{\textbf{Iacere}}, to throw.) What does \emph{-ivo} add? (\textbf{Not a new direction} — the sense of a thing placed there to be reached.)
\end{tcolorbox}
\section[(producir, suficiente, alimentar, el …]{(producir, suficiente, alimentar, el objetivo) — from cold}

\label{lesson:ES-R467-repaso-alimento}



\begin{quote}
Four words, no prompts. What the garden does, whether it is enough, who it keeps going, and the number written at the top of the sheet.
\end{quote}

\subsection*{You'll want to know: the four, cold}
\noindent
\begin{tabularx}{\linewidth}{@{}>{\raggedright\arraybackslash}X>{\raggedright\arraybackslash}X@{}}
\toprule
\textbf{English} & \textbf{Spanish} \\
\midrule
to produce & \textbf{producir} \\
enough & \textbf{suficiente} \\
to feed & \textbf{alimentar} \\
the target & \textbf{el objetivo} \\
\bottomrule
\end{tabularx}

\begin{grammarlens}[title={Grammar Lens: what this chapter actually cost you}]
Three of the four are assembled out of parts you already held.

\noindent
\begin{tabularx}{\linewidth}{@{}>{\raggedright\arraybackslash}X>{\raggedright\arraybackslash}X@{}}
\toprule
\textbf{word} & \textbf{where the parts came from} \\
\midrule
producir & \emph{reducir}'s verb class, \emph{ducere}'s prefix trick \\
objetivo & \emph{el objeto}'s throw, \emph{rotativo}'s \emph{-ivo} \\
\bottomrule
\end{tabularx}

Neither needed a new rule — only a new arrangement of rules you had.

\emph{Alimentar} borrowed \emph{-mento} from \emph{el medicamento}. Only \emph{suficiente} asked for something genuinely new, and even that was a distinction rather than a word: you already had \emph{bastante}.

That is what the earlier chapters bought. New words stop being new machinery and start being new arrangements of it, and each one costs less than the last.
\end{grammarlens}

\begin{cousinweb}
Two of the four hide a physical action.

\noindent
\begin{tabularx}{\linewidth}{@{}>{\raggedright\arraybackslash}X>{\raggedright\arraybackslash}X@{}}
\toprule
\textbf{word} & \textbf{what is underneath} \\
\midrule
suficiente & \emph{sub-} + \emph{facere} — \textbf{putting enough underneath} to hold it up \\
alimentar & \emph{alere} — \textbf{growing} a thing, with food as the means \\
\bottomrule
\end{tabularx}

\emph{Producir} leads forward; \emph{objetivo} is the thing placed in front to be reached. Four words off a garden shed, four Latin verbs of handling and growing underneath them — the pattern this book keeps finding.
\end{cousinweb}

\subsection*{Your turn}
Say these aloud:

\begin{itemize}
\raggedright
  \item we do not produce enough to feed everybody
  \item the month's target
  \item which of the four came with its irregular forms already learned
\end{itemize}

\begin{tcolorbox}[breakable,colback=teal!4,colframe=teal!35!black,title={Before you move on}]
To produce? (\textbf{Producir}.) Enough? (\textbf{Suficiente}.) To feed? (\textbf{Alimentar}.) The target? (\textbf{El objetivo}.)
\end{tcolorbox}
\section[Practice]{(el informe mensual del huerto) — read the figures, then judge them}

\label{lesson:ES-C467-sintesis-alimento}



\begin{quote}
Say \textbf{producir} and \textbf{el objetivo}. A sheet with both on it is asking one question: did the first reach the second?
\end{quote}

\begin{tcolorbox}[breakable,skin=enhanced,colback=orange!6,colframe=orange!45!black,boxrule=0.5pt,arc=1mm,left=6pt,right=6pt,top=4pt,bottom=4pt,fonttitle=\bfseries,title={Read this}]
\begin{quote}
\textbf{HUERTO COMUNITARIO — MARZO}
\end{quote}

\begin{quote}
\textbf{Objetivo} del mes: \textbf{alimentar} a 50 familias.
\end{quote}

\begin{quote}
\textbf{Producido}: 380 kg. Familias atendidas: 41.
\end{quote}

\begin{quote}
No fue \textbf{suficiente}. En abril vamos a plantar más patata (papa), que \textbf{produce} más kilos por metro cuadrado.
\end{quote}
\end{tcolorbox}

\begin{grammarlens}[title={Grammar Lens: the sheet answers its own question}]
Four numbers and one verdict, and the verdict is the word this chapter was for.

\noindent
\begin{tabularx}{\linewidth}{@{}>{\raggedright\arraybackslash}X>{\raggedright\arraybackslash}X@{}}
\toprule
\textbf{what the sheet gives} & \textbf{what it settles} \\
\midrule
50 against 41 & the target was missed \\
\emph{no fue suficiente} & and by how the garden judges it \\
\bottomrule
\end{tabularx}

It says \emph{suficiente}, not \emph{bastante}, and that is deliberate — though both would be understood. \emph{No fue bastante} is what you would say out loud. \emph{No fue suficiente} is what a sheet writes, because it names the requirement rather than the quantity, and the 380 kg is being measured against the 50 families.

Note \textbf{producido} too: the participle, used as a plain label for the amount. A sheet does not need a whole sentence to report a quantity.
\end{grammarlens}

\subsection*{Your turn}
Say in one sentence what the target was and whether it was met.

Now say what April is meant to change, and why the potato is the chosen crop.

\begin{tcolorbox}[breakable,colback=teal!4,colframe=teal!35!black,title={Before you move on}]
The month's target? (\textbf{Alimentar a 50 familias}.) Was it enough? (\textbf{No fue suficiente}.) Why that word on a sheet? (\textbf{Because it names the requirement}, where \emph{bastante} names a quantity.)
\end{tcolorbox}
